\section{The compact mean and Cameron map on time blocks}
\label{S:actual-cameron-mean}

Retain the nonzero original covariance in Hypothesis~\ref{H:homog}, its
amplitude, $a = \alpha_0$, and the real phases $E = H^{1/2}\times H^{-1/2}$ and
$F = E\oplus\R^2$. Write $\Lambda = (1-\Delta)^{1/2}$. The form $B_\Gamma$ and
constant $R_{\mathrm{sp}}$ are those of
Section~\ref{S:spatial-occupation-foundation}. Let $H_{\mathrm{sp}}$ be the real
covariance-null completion of $C_c^\infty(\R^d)$ with inner product $B_\Gamma$,
and $H_t$ the analogous completion of $C_c^\infty(I)$, $I = (0,1)$, with time
kernel $|s-t|^{-a}$. Let $H = H_I$ be the completion of smooth tests supported
inside $I\times\R^d$. It embeds isometrically as a closed subspace of the given
$H_W$. We retain the given $W$ restricted to it.

Fix $R > 0$ and real even $\eta_R \in C_c^\infty(\R^d)$ with $0 \le \eta_R \le
1$, equal to one on the radius-$R$ ball and supported in the closed radius-$2R$
ball. The physical mask is $m_R = \eta_R^2$. Every operator path below takes its
mask argument equal to $m_R$: at this fixed radius the base profile $\zeta(y) =
m_R(Ry)$ meets the earlier hypotheses. The root $\eta_R$ factors state tests; it
does not multiply an original noise class. All compactness conclusions here fix
$R$.

\subsection{Fractional products and the compact spatial map}

\begin{lemma}[Fractional covariance products]
  \label{L:cameron-fractional-products}
  For $\alpha/4 < s < 1/2$, products of real $H^s$ functions define elements of
  $H_{\mathrm{sp}}$, with
  \begin{equation}\label{E:cameron-fractional-product}
    \|fg\|_{H_{\mathrm{sp}}}
    \le C_s^2\sqrt{R_{\mathrm{sp}}}\|f\|_{H^s}\|g\|_{H^s}.
  \end{equation}
  Their Hilbert pairings equal their physical integrals. In particular,
  \begin{equation}\label{E:cameron-fractional-absolute}
    \int|f_1(x)f_2(x)f_3(x-y)f_4(x-y)|\ud x\,\Gamma(\ud y)
    \le C_s^4R_{\mathrm{sp}}\prod_{j=1}^4\|f_j\|_{H^s}.
  \end{equation}
\end{lemma}

\begin{proof}
  First approximate real $H^1$ factors by compact smooth ones. The absolute
  four-linear bound~\eqref{E:spatial-four-linear} and covariance PSD make their
  products Cauchy in $H_{\mathrm{sp}}$, with norm at most
  $\sqrt{R_{\mathrm{sp}}}N(f)N(g)$. Expanding product differences proves
  independence of approximation; physical continuity identifies all pairings.
  PSD for products also follows from $fg = ((f+g)^2-(f-g)^2)/4$.

  Use real radial Fourier bands $D_0 = 1_{\{|\xi|<1\}}$ and $D_j =
  1_{\{2^{j-1}\le|\xi|<2^j\}}$, $j \ge 1$. For $f \in H^s$ put $f_j = D_jf$.
  Each block is real and in $H^1$. Frequency support, Parseval and weighted
  Cauchy--Schwarz give
  \begin{equation}\label{E:cameron-dyadic-product}
    N(f_j)\le2^{j\alpha/4}\|f_j\|_2,\qquad
    \sum_{j\ge0}N(f_j)\le C_s\|f\|_{H^s}.
  \end{equation}
  Indeed, $\sum_j2^{-2j(s-\alpha/4)} < \infty$, while $\sum_j2^{2js}\|f_j\|_2^2
  \le 2^{2s}\|f\|_{H^s}^2$. The low block has factor one. No $L^p$ multiplier
  result is used. Also $\sum_j\|f_j\|_2 < \infty$ since $s > 0$. Thus,
  $\sum_j|f_j|$ lies in $L^2$ by Minkowski and monotone convergence, and the
  series is absolutely convergent a.e. with sum $f$.

  Apply the absolute $H^1$ four-linear estimate to each band quadruple. The sum
  of its absolute integrals is at most
  $R_{\mathrm{sp}}\prod_{\ell=1}^4\sum_jN((f_\ell)_j) < \infty$. Tonelli and the
  null-set transfer at both spatial positions in
  Lemma~\ref{L:spatial-physical-pair} identify the physical product series,
  proving~\eqref{E:cameron-fractional-absolute}. The double series
  $\sum_{j,k}f_jg_k$ is also absolutely convergent in $H_{\mathrm{sp}}$ by the
  $H^1$ product bound. Its partial sums are products of $H^1$ functions. The
  absolute quadruple calculation identifies their limiting physical pairings and
  gives~\eqref{E:cameron-fractional-product}. Compact smooth functions are dense
  in $H^s$: approximate in Fourier space by real Schwartz functions, then cut
  each off compactly in $H^1$ and hence in $H^s$. This also proves uniqueness of
  the bilinear extension. This proves Lemma~\ref{L:cameron-fractional-products}.
\end{proof}

\begin{proposition}[Compact spatial mean map]\label{P:cameron-spatial-map}
  The smooth-test map
  \[
    T_{\mathrm{sp}}q=iB_F(m_R(\Gamma*q))
  \]
  extends to a compact linear map $T_{\mathrm{sp}}:H_{\mathrm{sp}} \to
  \mathcal{K}(F)$, where $\mathcal{K}(F)$ has the operator norm. For $z =
  (z_1,z_2,c_1,c_2)$ and $w = (w_1,w_2,d_1,d_2)$ set $w' = \Lambda^{-1}w_2$. Its
  exact row identity is
  \begin{equation}\label{E:cameron-spatial-row}
    \langle T_{\mathrm{sp}}qz,w\rangle_F
    =\langle q,k(z,w)\rangle_{H_{\mathrm{sp}}},
    \qquad
    k(z,w)=(\eta_Rz_1+c_1\eta_R)(\eta_Rw').
  \end{equation}
\end{proposition}

\begin{proof}
  Fix $s$ with $\alpha/4 < s < 1/2$; this interval is nonempty because the
  original threshold implies $\alpha < 2$. For smooth compact $q$, $\Gamma*q$ is
  locally smooth by differentiation under the compactly restricted Radon
  integral. The masked potential is smooth compact. The velocity pairing in $E$
  has weight $\Lambda^{-1}$; hence physical Fubini
  gives~\eqref{E:cameron-spatial-row} on smooth states. Evenness of $\Gamma$
  changes $B_\Gamma(k,q)$ to $B_\Gamma(q,k)$. For general states the localized
  products are $L^1$, and $\Gamma*q$ is bounded on their compact support. $L^2$
  state approximation passes the physical integral to the limit, while
  Lemma~\ref{L:cameron-fractional-products} passes the $H_{\mathrm{sp}}$
  pairing. This proves the row.

  Multiplication by $\eta_R$ maps bounded $H^{1/2}$ sets to precompact $H^s$
  sets by Proposition~\ref{P:cutoff-singular-values}, with gain $1/2-s > 0$.
  This applies to $z_1$ and $w'$. The $c_1\eta_R$ family is finite dimensional
  and compact. The closures of both factor families for unit $z,w$ are compact,
  so their products have compact closure $\mathcal{L}$ in $H_{\mathrm{sp}}$ by
  the fractional product bound. The row identity gives $\|T_{\mathrm{sp}}q\| \le
  \sup_{k\in\mathcal{L}}\|k\|\|q\|$. Null covariance classes therefore map to
  zero.

  For smooth $q$, multiplication by its potential is compact $H^{1/2} \to
  H^{-1/2}$ by the same whole-space cutoff result with gain one. Its extra $c_1$
  column is finite rank, and the instantaneous $c_2$ column is zero. Thus, the
  smooth images lie in $\mathcal{K}(F)$, and so does the bounded extension by
  operator-norm closure.

  Choose a complete smooth-test orthonormal basis of $H_{\mathrm{sp}}$ by
  countable compact-support approximation and Gram--Schmidt, omitting null
  classes. Let $Q_n$ project onto its first $n$ terms. The row identity yields
  \[
    \|T_{\mathrm{sp}}(I-Q_n)\|
    \le\sup_{k\in\mathcal{L}}\|(I-Q_n)k\|_{H_{\mathrm{sp}}}
    \longrightarrow0.
  \]
  Strong projection convergence is uniform on compact $\mathcal{L}$ by a
  finite-net argument and $\|I-Q_n\| \le 1$. Each $T_{\mathrm{sp}}Q_n$ has
  finite-dimensional range; their operator-norm limit is a compact map. This is
  compactness of the map, in addition to compact values. This proves
  Proposition~\ref{P:cameron-spatial-map}.
\end{proof}

\subsection{The supported tensor and temporal mean}

\begin{lemma}[Supported original tensor space]\label{L:cameron-supported-tensor}
  Canonically $H = H_t\otimes H_{\mathrm{sp}}$ as real Hilbert spaces. It is
  separable and admits a complete orthonormal basis $e_j$ with smooth test
  representatives supported inside $I\times\R^d$. Projection $P_J$ onto its
  first $J$ elements is also an orthogonal projection on $H_W$ by the same
  finite formula.
\end{lemma}

\begin{proof}
  On separated smooth tests, absolute Fubini makes the original inner product
  exactly the tensor product after both null quotients. Finite sums are dense:
  approximate a smooth compact spacetime test uniformly by smooth time
  partitions times its spatial slices, on one common compact support inside
  $I\times\R^d$. A uniform error $\epsilon$ on spatial compact $A$ has squared
  original covariance norm at most
  \[
    \epsilon^2|A|\Gamma(A-A)
    \int_I\int_I|s-t|^{-a}\ud s\ud t<\infty.
  \]
  This also gives countable smooth-test density on exhausting compact sets, both
  here and in $H_{\mathrm{sp}}$. Gram--Schmidt skipping zeros preserves smooth
  compact support. The isometric embedding into $H_W$ has closed image by
  completeness; its finite orthonormal projection formula remains orthogonal
  there. It is not multiplication by a time indicator and gives no independence
  of disjoint time intervals. This proves
  Lemma~\ref{L:cameron-supported-tensor}.
\end{proof}

\begin{proposition}[The compact temporal mean, without trace class]
  \label{P:cameron-temporal-map}
  Let $K$ on $L^2(I)$ have kernel $|s-t|^{-a}$. It is compact, self-adjoint,
  strictly positive and injective. The map $\iota:f\mapsto K^{1/2}f$ extends
  isometrically from $H_t$ onto $L^2(I)$, and the temporal mean
  \begin{equation}\label{E:cameron-temporal-map}
    C_t=K^{1/2}\iota:H_t\longrightarrow L^r(I)
  \end{equation}
  is bounded and compact for every $2 < r < 2/a$. The displayed formula first
  has its $L^2(I)$ meaning; the same completed map is then shown to take values
  compactly in $L^r(I)$. The operator $K$ is not trace class and $C_t:H_t \to
  L^2(I)$ is not Hilbert--Schmidt.
\end{proposition}

\begin{proof}
  Schur bounds $K$ by its finite row and column integrals. The Gamma
  representation of $|s-t|^{-a}$ and zero extension of $f \in L^2(I)$ give
  strict positivity directly: the kernel $e^{-\lambda|v|}$ has quadratic form
  \[
    \int_\R\frac{2\lambda}{\lambda^2+\xi^2}
    |\widehat{f}(\xi)|^2\ud\xi>0
    \quad(f\ne0,\ \lambda>0)
  \]
  in the unitary Fourier convention. Absolute Fubini is justified by Schur
  applied to $|f(s)f(t)|$. Integrating with positive weight
  $\lambda^{a-1}/\Gamma_E(a)$ proves positivity and injectivity. Symmetry gives
  self-adjointness. The deterministic truncated kernel $K_\epsilon =
  K1_{\{|s-t|\ge\epsilon\}}$ is square integrable, so finite separated-kernel
  approximations make it compact. The Schur norm of $K-K_\epsilon$ is at most
  $2\epsilon^{1-a}/(1-a)$, proving compactness of $K$. These truncations are not
  used as Gaussian covariances. Its positive square root is compact by the
  compact spectral theorem.

  On tests $\|K^{1/2}f\|_2 = \|f\|_{H_t}$. Smooth-test $L^2$ density and $\ker K
  = 0$ imply that its range is dense, so the isometry extends onto $L^2(I)$. On
  tests $C_tf = Kf$, and~\eqref{E:cameron-temporal-map} first gives compactness
  into $L^2$.

  Extend $K_a(v) = |v|^{-a}1_{\{|v|\le1\}}$ to the line. It lies in $L^{r/2}$.
  Scalar Young convolution gives
  \[
    \|g\|_{H_t}^2\le\|K_a\|_{r/2}\|g\|_{r'}^2,\qquad
    r'=r/(r-1).
  \]
  To verify the needed Young inequality, take $u = r/2$, $v = r'$, so $u,v < r$
  and $1/u+1/v = 1+1/r$. Split the absolute convolution integrand into
  $(|f|^u|h|^v)^{1/r}$, $|f|^{1-u/r}$ and $|h|^{1-v/r}$. H\"older with exponents
  $r$, $(1/u-1/r)^{-1}$ and $(1/v-1/r)^{-1}$, followed by integration and
  Tonelli, gives $\|f*h\|_r \le \|f\|_u\|h\|_v$. Covariance Cauchy--Schwarz and
  $L^r$ duality now imply $\|Kf\|_r \le C_{a,r}\|f\|_{H_t}$ on tests. Their $Kf$
  are bounded by the integrable-kernel row, making duality legitimate before
  completion. Density extends the estimate, with $L^2$ value equal
  to~\eqref{E:cameron-temporal-map}. Choose $r < r_2 < 2/a$. Compactness into
  $L^2$ and boundedness into $L^{r_2}$ give compactness into $L^r$ by $\|u\|_r
  \le \|u\|_2^\tau\|u\|_{r_2}^{1-\tau}$ with $1/r = \tau/2+(1-\tau)/r_2$.

  Finally, the orthonormal cell indicators $N^{1/2}1_{(j/N,(j+1)/N)}$, $0 \le j
  < N$, have $K$ quadratic forms summing to $2N^a/((1-a)(2-a)) \to \infty$.
  Thus, positive $K$ has infinite trace, and $K^{1/2}\iota$ is not
  Hilbert--Schmidt into $L^2$. No instantaneous $L^2$-valued Gaussian has been
  introduced. This proves Proposition~\ref{P:cameron-temporal-map}.
\end{proof}

No-go~\ref{P:nogo-27} explains the upper endpoint of this compactness
range: temporal concentration makes the mean map unbounded above
$2/\alpha_0$ and noncompact at $2/\alpha_0$. The endpoint witness does not
refute boundedness there.


\begin{proposition}[The compact time-space mean]\label{P:cameron-full-mean}
  Fix one $2 < r < 2/a$ throughout this proposition. For finite tensors with
  orthonormal spatial labels define
  \[
    \mathcal{U}\left(\sum_j f_j\otimes q_j\right)(s)
    =\sum_j(Kf_j)(s)q_j.
  \]
  It extends to a bounded map $H \to L^r(I;H_{\mathrm{sp}})$ with dimension-free
  bound. The map
  \begin{equation}\label{E:cameron-full-mean}
    T\Phi(s)=T_{\mathrm{sp}}(\mathcal{U}\Phi(s))
    :H\longrightarrow L^r(I;\mathcal{K}(F))
  \end{equation}
  is compact. For smooth supported $\Phi$, it is the multiplier
  $iB_F(m_Rh_\Phi(s))$, where $h_\Phi = |\cdot|^{-a}*(\Gamma*\Phi)$. For the
  supported projections of Lemma~\ref{L:cameron-supported-tensor},
  \begin{equation}\label{E:cameron-compact-tail}
    \epsilon_J:=\|T(I-P_J)\|_{H\to L^r(\mathcal{K}(F))}
    \longrightarrow0.
  \end{equation}
\end{proposition}

\begin{proof}
  For a finite-dimensional spatial field $g$, positivity of the time kernel and
  scalar Young give
  \[
    \|g\|_{H_t\otimes H_{\mathrm{sp}}}^2
    \le\int_I\int_I|s-t|^{-a}\|g_s\|\|g_t\|\ud s\ud t
    \le C_{a,r}^2\|g\|_{L^{r'}(H_{\mathrm{sp}})}^2.
  \]
  Hilbert Cauchy--Schwarz and finite-dimensional $L^r$ duality give
  $\|\mathcal{U}\Phi\|_{L^r(H_{\mathrm{sp}})} \le C_{a,r}\|\Phi\|_H$. The
  duality can be checked by the pointwise direction
  $\|\mathcal{U}\Phi\|^{r-2}\mathcal{U}\Phi$ and smooth coefficient
  approximation in $L^{r'}$, controlled by the displayed estimate. No dimension
  factor occurs. Tensor density and completeness extend this to a strongly
  measurable Bochner field. This time-only map still contains spatial vectors in
  the covariance space. Spatial covariance is applied once, by
  $T_{\mathrm{sp}}$.

  The physical mean identity on separated tests follows
  from~\eqref{E:cameron-spatial-row}. For a smooth spacetime test use the
  separated compact-support approximations in the tensor proof. Their uniform
  errors give uniformly small masked mean errors, using the finite time row and
  locally finite spatial covariance on the relevant compact differences. The
  multiplier bound by $L^\infty+L^2$ then gives uniform operator convergence,
  identifying the $L^r$ mean in~\eqref{E:cameron-full-mean}. General $H$
  controls denote completed operator fields, without a pointwise-potential
  assertion.

  Let $Q_n$ be the spatial projections in the compact-map proof. The maps
  $T_n\Phi = T_{\mathrm{sp}}Q_n\mathcal{U}\Phi$ approximate $T$ in operator norm
  by the dimension-free bound and $\|T_{\mathrm{sp}}(I-Q_n)\| \to 0$. For fixed
  $n$, each spatial coordinate extraction $H \to H_t$ has norm at most one and
  its image under $C_t$ is precompact in $L^r$. Finite combination with the
  compact operators $T_{\mathrm{sp}}q_j$ makes $T_n$ compact. Hence, $T$ is
  compact. Scalar time compactness alone tensored with an infinite identity
  would not prove this.

  To prove~\eqref{E:cameron-compact-tail}, if it failed choose unit $v_J$ on a
  subsequence with $\|T(I-P_J)v_J\|$ bounded below. The vectors $(I-P_J)v_J$
  tend weakly to zero, since their pairings with fixed $h$ are bounded by
  $\|(I-P_J)h\|$. Compactness gives a norm-convergent image subsequence with
  weak limit zero, a contradiction. This proves
  Proposition~\ref{P:cameron-full-mean}.
\end{proof}

The radius remains fixed in this approximation. No-go~\ref{P:nogo-25}
shows that translated white-noise controls prevent these fixed smooth
projections from giving tails uniform over growing radii.


\subsection{The noise residual and its Cameron operator}

Choose integer $B \ge 1$ and $0 < \delta < \min(1,1/B)$. Put $N =
\lfloor1/(\delta B)\rfloor$, $s_m = m\delta B$, $\ell_m = B$ for $m < N$ and
$\ell_N = \delta^{-1}-NB \in [0,B)$. The intervals $I_m =
[s_m,s_m+\delta\ell_m]$ partition $[0,1]$ up to endpoints. With $Q_J = I-P_J$ on
$H_W$ and the explicit pullback~\eqref{E:raw-block-pullback}, define
\[
  Y_{m,J}(f)=W(Q_JT_{\delta,m}f),\qquad
  g_\delta=\delta^{(1-a)/2}.
\]
Fix parameters
\begin{equation}\label{E:cameron-block-parameters}
  \frac12<\theta<\min\left(1-\frac a2,\frac{3-a-\alpha}{2}\right),
  \qquad \frac1{1-\theta}<r<\frac2a,\qquad
  \frac1\theta\le p<2<r .
\end{equation}
These intervals are nonempty.

\begin{lemma}[Supported residual paths]\label{L:cameron-actual-residual}
  Gaussian transfer with mask $m_R$ and its $F$ extension constructs the
  centered Gaussian paths
  \[
    X_{m,J}(u)=g_\delta A_{m_R}(Y_{m,J})(\min(u,\ell_m)),
    \qquad 0\le u\le B.
  \]
  Their finite tuple lies in the separable space $\mathscr{B}_0^{N+1}$, where
  $\mathscr{B}_0 = h_0^\theta([0,B];\mathcal{K}(F))$. Every continuous linear
  statistic $\ell(X)$ is $W(k_\ell)$ for some $k_\ell \in H_I$. An empty last
  interval gives zero.
\end{lemma}

\begin{proof}
  Isometry of the pullback and orthogonality of $Q_J$ give
  $\operatorname{Var}Y_{m,J}(f) \le \|f\|_{H_W}^2$.
  Proposition~\ref{P:source-gaussian-transfer} therefore applies. Its finite
  sums retain both deterministic localizations, and their limits retain the true
  $F$ extension and constant columns. Choose $\theta < \beta < \beta_2$ below
  the upper threshold in~\eqref{E:cameron-block-parameters}; read the
  operator-path construction at exponent $\beta$ with auxiliary $\beta_2$.
  Constant stopping preserves compact values and $C^\beta$ bounds, hence gives
  paths in the little $C^\theta$ closure by polygonal approximation. Its
  separability is Lemma~\ref{L:raw-block-inputs}.

  As in Lemma~\ref{L:raw-block-inputs}, every observed row's slow pullback
  belongs to $H_I$: its time support is in $[0,1]$, and approximating its
  indicator inside $I$ converges in original covariance norm by dominated
  convergence against the integrable time kernel and locally finite $\Gamma$ on
  the common compact support. This includes endpoints. Schwartz $E$ states and
  both true constant columns supply the identities, then boundedness extends
  them to $F$. Let $P_H$ be the orthogonal projection onto closed $H_I$.
  Transfer for the noise $f\mapsto W(P_HQ_JT_{\delta,m}f)$ leaves each observed
  row unchanged, since its residual slow test already lies in $H_I$. Countable
  states and rational times identify the same operator paths. This
  projected-transfer series has coefficient tests in $H_I$. Applying $\ell$
  gives an $L^2$ limit in the closed isonormal image $W(H_I)$, proving the
  assertion. Original global-basis tests were not assumed individually
  supported. No time-support independence or Hilbert-valued noise is used. This
  proves Lemma~\ref{L:cameron-actual-residual}.
\end{proof}

\begin{proposition}[Cameron covariance and averaged bound]
  \label{P:cameron-actual-embedding}
  The covariance map $L_{\delta,J}\Phi = \E[XW(\Phi)]$ is
  \begin{equation}\label{E:cameron-actual-covariance}
    (L_{\delta,J}\Phi)_m(u)=\sqrt\delta
    \int_0^{\min(u,\ell_m)}
    S_F(-v)T(Q_J\Phi)(s_m+\delta v)S_F(v)\ud v ,
    \qquad \Phi\in H.
  \end{equation}
  This is a compact-operator Bochner integral. Define the exact norm and
  constants
  \begin{equation}\label{E:cameron-averaged-norm}
    \|x\|_{\delta,p}
    =\left(\delta\sum_{m=0}^N\|x_m\|_{C^\theta}^p\right)^{1/p},
    \quad G_B=\sup_{|v|\le B}\|S_F(v)\|,
    \quad C_B=G_B^2[B^{1-1/r}+B^{1-1/r-\theta}].
  \end{equation}
  Then the Gaussian Cameron embedding satisfies
  \begin{equation}\label{E:cameron-averaged-bound}
    \|L_{\delta,J}\Phi\|_{\delta,p}
    \le\sqrt\delta\,C_B(2/B)^{1/p-1/r}
    \epsilon_J\|\Phi\|_H .
  \end{equation}
  Convergence $\epsilon_J \to 0$ fixes $R$ and is independent of $\delta$.
\end{proposition}

\begin{proof}
  For smooth supported $\Phi$, $Q_J\Phi$ is another such test. The explicit
  pullback covariance~\eqref{E:raw-block-mean-jacobian} gives
  \[
    \E[Y_{m,J}(f)W(\Phi)]
    =\delta^{a/2}\int f(v,x)
    h_{Q_J\Phi}(s_m+\delta v,x)\ud x\ud v .
  \]
  Absolute Fubini also holds for signed compact tests. Multiplying by $g_\delta$
  gives $\sqrt\delta$. Explicit rows and~\eqref{E:cameron-spatial-row} give
  \eqref{E:cameron-actual-covariance}. The constant-velocity input acquires
  displacement $v$ under $S_F(v)$; the constant displacement remains one. Both
  columns are included.

  Both $S_F$ and its adjoint are strongly continuous on Hilbert $F$. In
  Bessel-unitary $E$ coordinates the Fourier matrix is
  \[
    \begin{pmatrix}
      \cos(v|\xi|)&\langle\xi\rangle\sin(v|\xi|)/|\xi|\\
      -|\xi|\sin(v|\xi|)/\langle\xi\rangle&\cos(v|\xi|)
    \end{pmatrix}.
  \]
  Its continuous zero quotient and common bound $C(1+B)$ apply also to its
  adjoint; dominated convergence proves the assertion. The constant block is
  finite dimensional. For compact $A$, conjugation $S_F(-v)AS_F(v)$ is norm
  continuous: strong convergence is uniform on compact images, and the right
  factor follows by adjoints. Joint norm continuity in $(v,A)$ and uniform
  bounds give Bochner measurability for the conjugated strongly measurable
  compact-operator field. No norm continuity of the free group itself is
  assumed.

  Write $M = TQ_J\Phi$. H\"older and the primitive increment estimate give
  \begin{equation}\label{E:cameron-local-primitive-bound}
    \|(L_{\delta,J}\Phi)_m\|_{C^\theta}
    \le\sqrt\delta\,C_B\delta^{-1/r}\|M\|_{L^r(I_m)}.
  \end{equation}
  Its primitive exponent is $1-1/r > \theta$. Simple compact-operator integrand
  approximation in $L^r$ puts the primitive in $\mathscr{B}_0$; constant
  stopping does not increase the H\"older norm. Affine pullback preserves null
  sets. For fixed $\delta,J$ the covariance map is bounded by Bochner
  Cauchy--Schwarz and finite tuple $L^2$ moments. The integral is bounded by the
  displayed estimate, so equality on smooth controls extends by density to all
  $H$. Dense states and rational times make it a full operator-path identity.

  Let $H_{\mathrm{seen}}$ be the closed span of the $k_\ell$ from
  Lemma~\ref{L:cameron-actual-residual}. Then $\ell(L_{\delta,J}\Phi) = \langle
  k_\ell,\Phi\rangle_H$ and $\ker L_{\delta,J} = H_{\mathrm{seen}}^\perp$. The
  first Wiener chaos definition of the Cameron space is therefore
  $L_{\delta,J}(H_{\mathrm{seen}})$ with the minimum $H$ preimage norm. Extra
  directions are removed by orthogonal projection, without an inverse
  covariance. A bound in the original $H$ norm is thus the embedding bound.

  Put $t_m = \|M\|_{L^r(I_m)} \ge 0$ and $n = N+1 \ge 2$. Finite power mean at
  exponent $r/p > 1$, applied to $t_m^p$, gives
  \[
    \left(\sum_m t_m^p\right)^{1/p}
    \le n^{1/p-1/r}\left(\sum_m t_m^r\right)^{1/r}.
  \]
  The true partition gives $\sum_m t_m^r = \|M\|_{L^r(I)}^r$, including an empty
  last interval, and $n \le 2/(\delta B)$. Hence, the exact $\delta$-weighted
  norm gives
  \[
    \|L_{\delta,J}\Phi\|_{\delta,p}
    \le\sqrt\delta\,C_B(\delta n)^{1/p-1/r}\|M\|_{L^r}
    \le\sqrt\delta\,C_B(2/B)^{1/p-1/r}
    \epsilon_J\|\Phi\|_H.
  \]
  The $B$ factor remains explicit; covariance and mask occur once. This proves
  Proposition~\ref{P:cameron-actual-embedding}.
\end{proof}

These conclusions fix the radius and make no uniform-radius projection tail
assertion. No concentration estimate or original moment upper bound follows from
this section alone.
